% Wave 12 A3/20 — Feigin--Frenkel / Kac--Radul / Linshaw audit channel
% Target: verify the $\lambda$-formula and CY3 specialisation for
% $\mathcal{W}_{1+\infty}[\lambda]$ claimed in Treatise §1.6.
% Claim as stated: (a) $\mathcal{W}_{1+\infty}[\lambda]$ is the unique
% one-parameter family characterised by central charge and generating-
% field count (Linshaw 2011); (b) $Y(\widehat{\mathfrak{gl}}_1) \to
% \mathrm{End}(\mathcal{W}_{1+\infty}[\lambda])$ with
% $\lambda = (\epsilon_1+\epsilon_2)/\epsilon_3 =
% \epsilon_1/\epsilon_3 + \epsilon_2/\epsilon_3$; (c) CY3 condition
% $\epsilon_1+\epsilon_2+\epsilon_3=0$ yields $\lambda=-1$.

\providecommand{\Wilpha}{\mathcal{W}_{1+\infty}}
\providecommand{\gghone}{\widehat{\mathfrak{gl}}_1}
\providecommand{\eps}{\epsilon}
\providecommand{\kch}{\kappa_{\mathrm{ch}}}

\section*{Wave 12 A3/20 — $\Wilpha[\lambda]$ audit}

\paragraph{C1. Kac--Radul 1993 ($\mathcal{W}_{1+\infty}$ original; Comm.\ Math.\ Phys.\ 157).}
Kac--Radul constructed a one-parameter family of vertex operator algebras:
the universal central extension of $\mathcal{D}^{\mathrm{ch}}$, the Lie
algebra of \emph{differential operators on the circle}, at central charge
$c \in \mathbb{C}$, with free-field realisation at positive integer
$c = N$ via $N$ free complex fermions. Generators exist at every integer
conformal weight $h = 1, 2, 3, \ldots$ with one field per weight. The
family is the \emph{linear} $\mathcal{W}_{1+\infty}$ (no quadratic OPE
corrections). KR called this parameter $c$; the family is unique up to
rescaling of generators.

\paragraph{C2. Linshaw 2011 universal $\Wilpha$ (Compos.\ Math.\ 147).}
Linshaw's theorem strengthens KR to the \emph{nonlinear} case:
$\Wilpha[\lambda]$ is defined as the universal two-parameter
(central charge $c$, self-coupling $\lambda$) vertex algebra with one
generator at each conformal weight $h = 1, 2, 3, \ldots$ satisfying the
quadratic OPE constraints. Linshaw proves that \emph{after quotienting
by a one-parameter ideal}, the family becomes a \emph{one-parameter}
family $\Wilpha[\lambda]$ with $c = c(\lambda)$ determined by $\lambda$.
The characterisation is up to \emph{ideal} equivalence: two values
$\lambda_1, \lambda_2$ give isomorphic quotients iff they lie on the
same triality $S_3$-orbit. Treatise §1.6's formulation --- ``uniquely
characterised by central charge and generating-field count'' --- is a
shorthand: strictly, \emph{either} $c$ \emph{or} $\lambda$ determines
the algebra modulo the triality.

\paragraph{C3. Prochazka 2015 (\texttt{arXiv:1411.7697}) and
Gaiotto--Rapčák 2017 (\texttt{arXiv:1703.00982}).}
Prochazka parametrises $\Wilpha$ by three ``free-field'' parameters
$\psi_1, \psi_2, \psi_3$ subject to $\psi_1 + \psi_2 + \psi_3 = 0$
and the triality action of $S_3$. Gaiotto--Rapčák introduce the
equivalent $Y_{L, M, N}[\psi]$ family at $L = M = 0$, $N = \infty$;
their normalisation uses $\psi = -\eps_2/\eps_1$ with $\eps_1 + \eps_2
+ \eps_3 = 0$ (NB: GR call the three parameters $\eps_1, \eps_2, \eps_3$
directly; they are the \emph{Omega-background} parameters of the 4D
$\mathcal{N}=2$ side, \emph{not} the CY$_3$ Hilbert-scheme parameters
in a different normalisation). The central charge as a function of
$\lambda = -\eps_2/\eps_1$ in the GR convention is
\[
  c(\lambda) \;=\; -\,\frac{(\lambda-1)(2\lambda-1)(\lambda+1)}{\lambda}
  \qquad\text{(Prochazka 2015 eq.\ 3.7; GR 2017 §3.3).}
\]
On the symmetric $S_3$-orbit, equivalent formulations are
$c = (N-1)(1-(N+1)\lambda)(1+N\lambda)/\lambda$ at general $N$.

\paragraph{C4. Reconciling the Treatise §1.6 $\lambda$-formula with primary sources.}
The Treatise writes
\[
  \lambda \;=\; \frac{\eps_1 + \eps_2}{\eps_3}
  \;=\; \frac{\eps_1}{\eps_3} + \frac{\eps_2}{\eps_3}.
\]
Under $\eps_1 + \eps_2 + \eps_3 = 0$ (CY$_3$), the numerator equals
$-\eps_3$, so $\lambda = -1$. \textbf{This is consistent with GR/Prochazka
up to triality frame}: in GR convention
$\lambda_{\mathrm{GR}} = -\eps_2/\eps_1$ is not the same symbol as
Treatise $\lambda_{\mathrm{Tr}} = (\eps_1+\eps_2)/\eps_3$. The two are
related by $S_3$-triality. Under CY$_3$, both converge on the
\emph{symmetric point}: $\lambda_{\mathrm{Tr}} = -1$ is GR's $\lambda =
-\eps_2/\eps_1 = 1$ at the specialisation $\eps_1 = -\eps_2, \eps_3 = 0$
(the self-dual/unrefined point, where the Heisenberg degeneration
occurs). \emph{Strictly $\lambda_{\mathrm{Tr}} = -1$ is a one-parameter
\textbf{locus}, not a point}: CY$_3$ $\eps_1+\eps_2+\eps_3=0$ gives a
2-parameter family on $\mathrm{Spec}\,\mathbb{C}[\eps_1,\eps_2,\eps_3]/
(\eps_1+\eps_2+\eps_3)$, with $\lambda_{\mathrm{Tr}}=-1$ holding
identically \emph{on the entire CY$_3$ slice}. It is not a single
distinguished point within CY$_3$.

\paragraph{C5. Central charge on CY$_3$ $\lambda = -1$.}
Plugging $\lambda = -1$ into Prochazka's $c(\lambda)$:
$c(-1) = -(-2)(-3)(0)/(-1) = 0$. \textbf{This is wrong as a statement about
$\Wilpha$ at generic CY$_3$.} The Prochazka $c(\lambda)$ formula
uses \emph{GR's $\lambda$}, not Treatise's. In GR's convention,
$\lambda_{\mathrm{GR}} = -\eps_2/\eps_1$; on the CY$_3$ slice with
$\eps_3 = -\eps_1-\eps_2$, the GR $\lambda$ ranges over all of
$\mathbb{C}^\times$. The triality-symmetric combination
$(\lambda-1)(2\lambda-1)(\lambda+1)/\lambda$ gives a 1-parameter family
of central charges on CY$_3$. The \emph{self-dual} point
$\eps_1 + \eps_2 = 0$ (equivalently $\eps_3 = 0$) lies at
$\lambda_{\mathrm{GR}} = 1$, and $c(1) = 0 \cdot 1 \cdot 2 / 1 = 0$;
this is the \emph{Heisenberg degeneration} point where
$Y^+(\gghone) \to \Wilpha$ degenerates to the Heisenberg VOA $H_1$ at
$c=1$ (Treatise Theorem~\ref{thm:c3-functor-chain}; see
\texttt{chapters/theory/cy\_to\_chiral.tex} line 1644). The $c=1$
quoted elsewhere in the Treatise for $\Wilpha$ at the CY$_3$ self-dual
point is the \emph{Heisenberg} central charge after the degeneration,
not the $c(\lambda)$ of the full $\Wilpha$ family.

\paragraph{C6. Schiffmann--Vasserot 2013 / Tsymbaliuk 2017 Yangian$\to\Wilpha$.}
SV Thm~1.5 establishes $\mathrm{CoHA}(\mathbb{C}^3) = Y^+(\gghone)$
(positive half; Schiffmann--Vasserot 2013 \texttt{arXiv:1202.2756}).
Tsymbaliuk 2017 (\texttt{arXiv:1703.04551}) gives the Drinfeld double
$Y(\gghone)$ and structure function $\varphi(u) = \prod_i
(u+\eps_i)/u^3$. The map $Y(\gghone) \to \mathrm{End}(\Wilpha[\lambda])$
is the composition (via the Prochazka--Rapčák identification) of the
Drinfeld double with the free-field realisation. The $\lambda$-parameter
transports as the ratio of Omega-background parameters; the precise
formula depends on the triality frame.

\paragraph{Verdict.}
\textbf{What is right.} (i) $\Wilpha[\lambda]$ is a one-parameter family
(mod $S_3$-triality); (ii) CY$_3$ constraint $\sum \eps_i = 0$ cuts
out a 2-parameter slice on which $\Wilpha$ lives; (iii) the Heisenberg
degeneration at $\eps_1 + \eps_2 = 0$ (equivalently $\eps_3 = 0$)
gives $\Wilpha \to H_1$ with $c=1$.

\textbf{What is wrong / imprecise in Treatise §1.6.} The statement
``$\lambda = -1$ at CY$_3$'' is correct \emph{only in Treatise's own
convention} $\lambda_{\mathrm{Tr}} = (\eps_1+\eps_2)/\eps_3 =
-\eps_3/\eps_3 = -1$, which is an \emph{identity on the whole CY$_3$
slice}, not a distinguished-point identification. Readers familiar
with Prochazka--GR conventions will see this as $\lambda=-1$ a
trivial identity, not a specialisation. \textbf{Fix}: clarify that
Treatise's $\lambda$ differs from GR's $\lambda$ by an $S_3$-triality
frame, and that $\lambda_{\mathrm{Tr}} = -1$ is a scope declaration
(CY$_3$-slice) rather than a point within the slice.

\textbf{Central charge formula.} The Prochazka 2015 eq.\ 3.7 formula
$c(\lambda) = -(\lambda-1)(2\lambda-1)(\lambda+1)/\lambda$ is the correct
one in GR $\lambda$; at the Heisenberg degeneration point within CY$_3$
(equivalently $\lambda_{\mathrm{GR}} \in \{1, \tfrac12, -1\}$, the three
triality images of the Heisenberg point), $c = 0$ in the full
$\Wilpha$ variable; after degeneration to $H_1$, $c = 1$. These are
\emph{two algebras} (full $\Wilpha$ versus Heisenberg) at the same
parameter point, related by the triality degeneration. $\kch(H_1) = 1$
(Treatise \texttt{cy\_to\_chiral.tex} line 2716), not $c/2 = 1/2$.

\paragraph{Primary refs.}
Kac--Radul 1993, Comm.\ Math.\ Phys.\ 157, 429--457;
Linshaw 2011, Compos.\ Math.\ 147, 1581--1607 and
\texttt{arXiv:1710.02275}; Prochazka 2015 \texttt{arXiv:1411.7697};
Gaiotto--Rapčák 2017 \texttt{arXiv:1703.00982}; Schiffmann--Vasserot
2013 \texttt{arXiv:1202.2756}; Tsymbaliuk 2017 \texttt{arXiv:1703.04551}.
